\section{Avvio e utilizzo}\label{sec:usage}

L'intera architettura, comprendente tutti i servizi, con il codice nella
cartella \textit{src/}, può essere avviata e gestita tramite i comandi di
\textit{docker compose}; per comodità, sono presenti due script \textit{bash}
(per sistemi Linux). \textit{start.sh} lancia l'infrastruttura:
\begin{lstlisting}[language=Bash]
	# dalla cartella di base del repository
	./src/start.sh
\end{lstlisting}
\textit{down.sh} interrompe i servizi:
\begin{lstlisting}[language=Bash]
	./src/down.sh
\end{lstlisting}


Prima, però, è necessario scaricare la mappa (con dati di OpenStreetMap)
dell'Italia nord-orientale, disponibile al link su geofabrik:
\url{https://download.geofabrik.de/europe/italy/nord-est-latest.osm.pbf}.\\
Il file dovrà essere piazzato in \textit{service-graphhopper/data/}, con nome
del file\\
\mbox{\textit{geofabrik\_italy-nord-est-260517.osm.pbf}}.\\
Da Linux, si può eseguire direttamente, per entrambi i passaggi, il comando:
\begin{lstlisting}[language=Bash]
	# dalla cartella di base del repository
	wget -O service-graphhopper/data/geofabrik_italy-nord-est-260517.osm.pbf \
	https://download.geofabrik.de/europe/italy/nord-est-latest.osm.pbf
\end{lstlisting}


Il flusso delle attività parte sempre dall'utente, che in uno scenario reale
sarebbe un'entità fisica e potrebbe interagire con ACMEMobility tramite
un'applicazione per smartphone. Dunque, abbiamo realizzato uno script in
\textit{python} (\textit{src/user/user.py}) che simula tutte le possibili azioni
che possa compiere, in base agli argomenti utilizzati. Seguono degli esempi per
ogni utilizzo. Usage:
\begin{lstlisting}[language=Bash]
	# dalla cartella di base del repository
	python src/user/user.py
	# Usage: python user.py <action> <user_id> <vehicle_id> [station_id]
	# Actions: scan, reserve, cancel, scan_reserved, park, lock, lock_assistance, steal
\end{lstlisting}

Scansione veicolo (con id \textit{V-01}) per uso immediato, da parte dell'utente
con id \textit{USER-001}:
\begin{lstlisting}[language=Bash]
	python src/user/user.py scan USER-001 V-01
\end{lstlisting}

Prenotazione veicolo per dopo:
\begin{lstlisting}[language=Bash]
	python src/user/user.py reserve USER-001 V-01
\end{lstlisting}

Cancellazione della prenotazione precedente:
\begin{lstlisting}[language=Bash]
	python src/user/user.py cancel USER-001 V-01
\end{lstlisting}

Scansione del veicolo precedentemente prenotato:
\begin{lstlisting}[language=Bash]
	python src/user/user.py scan_reserved USER-001 V-01
\end{lstlisting}

Parcheggio del veicolo (simulazione dell'inserimento del veicolo nel dock della
stazione con id \mbox{\textit{STATION-BOLOGNA-01}}):
\begin{lstlisting}[language=Bash]
	python src/user/user.py park USER-001 V-01 STATION-bologna-01
\end{lstlisting}

Richiesta di fine corsa e, dunque, blocco del veicolo precedentemente
parcheggiato:
\begin{lstlisting}[language=Bash]
	python src/user/user.py lock USER-001 V-01 STATION-bologna-01
\end{lstlisting}

Richiesta di assistenza a fine corsa, in caso di problemi con il blocco:
\begin{lstlisting}[language=Bash]
	python src/user/user.py lock_assistance USER-001 V-01
\end{lstlisting}

Simulazione furto di un veicolo:
\begin{lstlisting}[language=Bash]
	python src/user/user.py steal USER-001 V-01
\end{lstlisting}

Invia un'istruzione a Camunda per eseguire immediatamente il task di timeout
per il tempo di cancellazione di una prenotazione (utile per testare
l'applicazione della penalità):
\begin{lstlisting}[language=Bash]
	python3 src/user/user.py force_timeout USER-001 V-01
\end{lstlisting}


Una volta avviata l'infrastruttura, è possibile accedere alla visualizzazione di
Camunda al link \\
\mbox{\url{http://localhost:8080/camunda/app/cockpit/}}, con username
\textit{demo} e password \textit{demo}.\\
Al link \mbox{\url{http://localhost:4000/map/}}, invece, è possibile
visualizzare una mappa di Bologna, con dei segnali per le stazioni ed eventuali
veicoli in movimento.\\
I veicoli utilizzano coordinate reali e ricevono dei percorsi da un'istanza,
eseguita in locale, di \textit{GraphHopper} --- che espone anche un'interfaccia
web al link \mbox{\url{http://localhost:8989/}}, non particolarmente rilevante
per il progetto.\\
\\
Sono inoltre disponibili altri endpoint, esposti dai vari servizi per
debugging. Seguono quelli più utili per il debugging.\\

Il servizio di ACMEMobility espone il suo database:  
\begin{lstlisting}[language=Bash]
	curl localhost/db
\end{lstlisting}

Ogni stazione espone il suo stato:
\begin{lstlisting}[language=Bash]
	curl localhost:5001/station
	curl localhost:5002/station
	curl localhost:5003/station
	curl localhost:5004/station
	curl localhost:5005/station
\end{lstlisting}
